% =========================================================
% CHAPITRE 4 : SPRINT 2
% =========================================================

\clearpage

\refstepcounter{chapter}

\setcounter{section}{0}
\setcounter{subsection}{0}
\setcounter{subsubsection}{0}

\phantomsection
\addcontentsline{toc}{chapter}
{Chapitre \thechapter : Sprint 2 -- Développement applicatif et moteur d'entraînement MLOps}

\markboth{Chapitre \thechapter}{Chapitre \thechapter}

\thispagestyle{empty}


% =========================================================
% PAGE D'OUVERTURE + PLAN
% =========================================================

\vspace*{0.8cm}

\begin{center}

\rule{0.30\textwidth}{1.2pt}
\hspace{0.5cm}
\textit{Chapitre \thechapter}
\hspace{0.5cm}
\rule{0.30\textwidth}{1.2pt}

\vspace{0.9cm}

\rule{\textwidth}{0.8pt}

\vspace{0.35cm}

{\fontsize{21}{24}\selectfont
\textsc{
\parbox{0.95\textwidth}{
\centering
Chapitre \thechapter : Sprint 2 -- Développement applicatif\\
et moteur d'entraînement MLOps
}
}
}

\vspace{0.35cm}

\rule{\textwidth}{0.8pt}

\end{center}


\vspace{1.4cm}

{\Large\bfseries Plan}

\vspace{0.65cm}

\begin{center}

\renewcommand{\arraystretch}{1.35}

\begin{tabular}{p{0.8cm}p{10.7cm}r}

\textbf{1}
&
\textbf{Introduction}
&
\pageref{sec:ch4-introduction}
\\[0.22cm]

\textbf{2}
&
\textbf{Objectifs et travaux du Sprint 2}
&
\pageref{sec:sprint2-objectifs}
\\[0.22cm]

\textbf{3}
&
\textbf{Développement de la partie applicative}
&
\pageref{sec:sprint2-developpement-applicatif}
\\[0.22cm]

\textbf{4}
&
\textbf{Gestion des projets, fichiers et types de données}
&
\pageref{sec:sprint2-projets-donnees}
\\[0.22cm]

\textbf{5}
&
\textbf{Gestion des métriques et du score global}
&
\pageref{sec:sprint2-evaluation}
\\[0.22cm]

\textbf{6}
&
\textbf{Mise en place du moteur d'entraînement}
&
\pageref{sec:sprint2-training-engine}
\\[0.22cm]

\textbf{7}
&
\textbf{Suivi des expériences, notifications et résultats}
&
\pageref{sec:sprint2-tracking-results}
\\[0.22cm]

\textbf{8}
&
\textbf{Diagrammes de séquence UML}
&
\pageref{sec:sprint2-uml}
\\[0.22cm]

\textbf{9}
&
\textbf{Bilan du Sprint 2}
&
\pageref{sec:sprint2-bilan}
\\[0.22cm]

\textbf{10}
&
\textbf{Conclusion}
&
\pageref{sec:ch4-conclusion}
\\[0.22cm]

\end{tabular}

\renewcommand{\arraystretch}{1}

\end{center}

\clearpage


% =========================================================
% 4.1 INTRODUCTION
% =========================================================

\section{Introduction}
\label{sec:ch4-introduction}

Après la préparation de l'environnement technique réalisée pendant le
premier sprint, le travail s'est concentré sur le développement des
premières fonctionnalités réellement utilisables de la plateforme.

Ce deuxième sprint marque le passage d'un environnement simplement
opérationnel vers une application capable de gérer les utilisateurs,
les projets Machine Learning et les fichiers nécessaires aux
entraînements.

Il couvre également la prise en charge de plusieurs types de données,
la gestion des métriques, le calcul du score global, la mise en place
des deux modes d'entraînement ainsi que le suivi des expériences avec
MLflow.

Le développement a donc été mené sur deux axes complémentaires. Le
premier concerne la partie applicative, avec l'évolution du frontend
Angular et du backend Spring Boot. Le second concerne le moteur MLOps,
avec les mécanismes nécessaires à la préparation, à l'exécution et au
suivi des entraînements.

L'objectif de ce sprint était d'obtenir un premier parcours cohérent :
s'authentifier, créer un projet, importer les fichiers et les données,
configurer l'évaluation, lancer un entraînement puis consulter son état
et ses résultats.


% =========================================================
% 4.2 OBJECTIFS ET TRAVAUX
% =========================================================

\section{Objectifs et travaux du Sprint 2}
\label{sec:sprint2-objectifs}

Le Sprint 2 s'est déroulé du \textbf{19 mars au 16 mai}. Les principaux
travaux réalisés pendant cette période sont résumés dans le tableau
\ref{tab:travaux-sprint2}.

\begin{table}[H]

\centering
\small

\renewcommand{\arraystretch}{1.16}
\setlength{\tabcolsep}{5pt}

\begin{tabular}{|p{4cm}|p{9.2cm}|}

\hline

\rowcolor{cyan!20}
\textbf{Élément}
&
\textbf{Rôle dans le Sprint 2}
\\
\hline

\textbf{Authentification et profils}
&
Permettre l'accès à la plateforme, la gestion du compte utilisateur
et la séparation des droits entre utilisateur et administrateur.
\\
\hline

\rowcolor{cyan!7}
\textbf{Gestion des projets et fichiers}
&
Créer les projets Machine Learning et importer les datasets, scripts
Python, fichiers de dépendances et Dockerfiles.
\\
\hline

\textbf{Types de données}
&
Prendre en charge les données tabulaires, textuelles et images afin
de couvrir plusieurs cas d'usage Machine Learning.
\\
\hline

\rowcolor{cyan!7}
\textbf{Gestion des métriques}
&
Adapter les métriques utilisées au type de problème :
classification ou régression.
\\
\hline

\textbf{Score global}
&
Combiner les métriques sélectionnées selon leur ordre de priorité
afin de disposer d'un indicateur commun pour comparer les expériences.
\\
\hline

\rowcolor{cyan!7}
\textbf{Mode standard}
&
Permettre un entraînement rapide à partir d'une image Docker fixe
déjà préparée.
\\
\hline

\textbf{Mode Docker personnalisé}
&
Construire un environnement spécifique au projet à partir du
Dockerfile et des dépendances fournis par l'utilisateur.
\\
\hline

\rowcolor{cyan!7}
\textbf{MLflow}
&
Conserver le suivi des runs, des métriques et des artefacts produits
pendant les entraînements.
\\
\hline

\textbf{Notifications et résultats}
&
Informer l'utilisateur du succès ou de l'échec d'un run et lui
permettre de consulter les résultats depuis la plateforme.
\\
\hline

\end{tabular}

\caption{Principaux travaux réalisés lors du Sprint 2}
\label{tab:travaux-sprint2}

\end{table}


% =========================================================
% 4.3 DÉVELOPPEMENT APPLICATIF
% =========================================================

\section{Développement de la partie applicative}
\label{sec:sprint2-developpement-applicatif}

Le développement applicatif a évolué en parallèle du moteur MLOps.
Le frontend devait proposer un parcours simple à l'utilisateur, tandis
que le backend devait centraliser la logique métier et assurer les
échanges avec la base de données, Kubernetes et les services MLOps.


% =========================================================
\subsection{Authentification et gestion des profils}
\label{subsec:sprint2-auth}

La première fonctionnalité développée a été l'authentification. Elle
repose sur un système d'inscription, de connexion et de gestion des
rôles, avec une vérification par email.

Lors de l'inscription, l'utilisateur fournit son nom, son email et
un mot de passe. Un email de vérification est automatiquement envoyé
avant de pouvoir accéder à la plateforme. La connexion s'effectue
ensuite via un token JWT qui sécurise les échanges entre le frontend
et le backend.

Un mécanisme de récupération de mot de passe a également été mis en
place. L'utilisateur peut demander la réinitialisation de son mot de
passe depuis la page de connexion. Un email contenant un lien
temporaire lui est envoyé, lui permettant de définir un nouveau mot
de passe.

Côté frontend, trois guards contrôlent l'accès aux pages :
\texttt{guestGuard} pour les pages d'inscription et de connexion
(accessibles uniquement aux utilisateurs non connectés),
\texttt{authGuard} pour les pages protégées, et
\texttt{adminGuard} pour les pages réservées à l'administration.

Deux acteurs principaux ont été définis. L'\textbf{utilisateur
standard} dispose de toutes les fonctionnalités de la plateforme :
création de projets, importation de fichiers, lancement d'entraînements,
suivi des expériences, déploiement et monitoring. L'\textbf{administrateur}
peut en outre gérer les comptes utilisateurs et consulter une vue
globale sur l'état de la plateforme.

Le diagramme de séquence \ref{fig:uml-auth-sequence} illustre le flux
complet d'authentification, depuis l'inscription jusqu'à la connexion.
L'utilisateur soumet ses informations via le frontend Angular, le
backend Spring Boot valide les données, insère le compte dans MySQL
et déclenche l'envoi d'un email de vérification. Une fois le compte
confirmé, l'utilisateur se connecte et reçoit un token JWT qui lui
donne accès aux fonctionnalités protégées de la plateforme.

\begin{figure}[H]
    \centering

    \includegraphics[
        width=0.92\textwidth,
        keepaspectratio
    ]{images/uml_auth_sequence.png}

    \caption{Diagramme de séquence -- Authentification}
    \label{fig:uml-auth-sequence}
\end{figure}

\begin{figure}[H]
    \centering

    \includegraphics[
        width=0.76\textwidth,
        keepaspectratio
    ]{images/sprint2_authentication.png}

    \caption{Interface d'authentification de la plateforme}
    \label{fig:sprint2-authentication}
\end{figure}

La figure \ref{fig:sprint2-authentication} présente les pages
d'inscription et de connexion. Après authentification, l'utilisateur
est redirigé vers le tableau de bord et retrouve son espace personnel
ainsi que les fonctionnalités associées à son compte.

\begin{figure}[H]
    \centering

    \includegraphics[
        width=0.78\textwidth,
        keepaspectratio
    ]{images/sprint2_profile.png}

    \caption{Interface de gestion du profil utilisateur}
    \label{fig:sprint2-profile}
\end{figure}

La figure \ref{fig:sprint2-profile} montre la page de profil. Elle
permet à l'utilisateur de consulter et de modifier ses informations
personnelles ainsi que de gérer les paramètres de son compte.


% =========================================================
\subsection{Frontend Angular et backend Spring Boot}
\label{subsec:sprint2-front-back}

Le frontend Angular a été enrichi avec les interfaces nécessaires à la
création des projets, à l'importation des fichiers, au lancement des
expériences et à la consultation des résultats.

En parallèle, le backend Spring Boot a été étendu avec les API REST et
les services nécessaires à la gestion de ces opérations. Il joue le
rôle d'intermédiaire entre l'interface, MySQL et les différents
composants du moteur MLOps.

Swagger a également été utilisé pendant le développement afin de tester
les endpoints de manière indépendante avant leur intégration complète
dans l'interface.


% =========================================================
% 4.4 PROJETS, FICHIERS ET DONNÉES
% =========================================================

\section{Gestion des projets, fichiers et types de données}
\label{sec:sprint2-projets-donnees}


% =========================================================
\subsection{Gestion des projets et des fichiers}
\label{subsec:sprint2-projets-fichiers}

La gestion des projets constitue le point de départ du workflow
d'entraînement. Un projet regroupe les informations et les fichiers
nécessaires à l'exécution d'une expérience.

Pour le mode standard, l'utilisateur fournit principalement son script
Python et son dataset. Pour le mode Docker personnalisé, il peut
également fournir un fichier \textit{requirements.txt} ainsi qu'un
Dockerfile afin de définir son propre environnement d'exécution.

Cette organisation permet de conserver un lien clair entre le projet,
les fichiers utilisés, le run lancé et les résultats obtenus.

\begin{figure}[H]
    \centering

    \includegraphics[
        width=0.80\textwidth,
        keepaspectratio
    ]{images/sprint2_projects.png}

    \caption{Interface de gestion des projets Machine Learning}
    \label{fig:sprint2-projects}
\end{figure}

La gestion des fichiers complète ensuite cette étape en permettant
d'associer au projet les ressources nécessaires à l'entraînement.

\begin{figure}[H]
    \centering

    \includegraphics[
        width=0.80\textwidth,
        keepaspectratio
    ]{images/sprint2_upload_files.png}

    \caption{Importation des fichiers nécessaires à l'entraînement}
    \label{fig:sprint2-upload-files}
\end{figure}

% =========================================================
\subsection{Prise en charge de différents types de données}
\label{subsec:sprint2-data-types}

La plateforme a été conçue pour prendre en charge plusieurs types de
données afin de couvrir différents cas d'usage Machine Learning.

Lors de la création d'un projet, c'est \textbf{l'utilisateur qui choisit
le type de données} qu'il souhaite utiliser. Trois catégories principales
sont supportées :

\begin{itemize}
    \item les \textbf{données tabulaires}, principalement sous forme de fichiers CSV ;
    \item les \textbf{données textuelles}, utilisées pour les projets NLP ;
    \item les \textbf{données images}, utilisées notamment pour la classification d'images.
\end{itemize}

Ce choix est associé au projet et permet à la plateforme de gérer
correctement les données utilisées lors des entraînements.


% =========================================================
% 4.5 MÉTRIQUES ET SCORE GLOBAL
% =========================================================

\section{Gestion des métriques et du score global}
\label{sec:sprint2-evaluation}


% =========================================================
\subsection{Métriques selon le type de problème}
\label{subsec:sprint2-metrics}

Les performances d'un modèle ne s'évaluent pas de la même manière
selon qu'il s'agit d'un problème de classification ou de régression.

La plateforme adapte donc les métriques proposées au type de projet.

Pour la classification, les principales métriques utilisées sont
l'Accuracy, la Precision, le Recall et le F1-score.

Pour la régression, l'évaluation repose principalement sur le MAE,
le MSE, le RMSE et le coefficient $R^2$.

Chaque métrique est associée à une \textbf{direction} qui indique
si une valeur élevée ou basse est préférable : \textbf{max} pour
les métriques où une valeur élevée est souhaitée (Accuracy, Precision,
Recall, F1, $R^2$) et \textbf{min} pour les métriques où une valeur
basse est préférable (RMSE, MAE, MSE, loss).

\begin{table}[H]

\centering
\small

\renewcommand{\arraystretch}{1.15}

\begin{tabular}{|p{4cm}|p{6.5cm}|p{2cm}|}

\hline

\rowcolor{cyan!20}
\textbf{Type de problème}
&
\textbf{Métriques principales}
&
\textbf{Direction}
\\
\hline

\textbf{Classification}
&
Accuracy, Precision, Recall, F1-score
&
max
\\
\hline

\rowcolor{cyan!7}
\textbf{Régression}
&
MAE, MSE, RMSE, $R^2$
&
min (sauf $R^2$ : max)
\\
\hline

\end{tabular}

\caption{Métriques utilisées selon le type de problème}
\label{tab:sprint2-metrics}

\end{table}

Cette distinction permet de conserver une évaluation cohérente avec la
nature du modèle et d'éviter de proposer des métriques qui ne
correspondent pas au cas d'usage.


% =========================================================
\subsection{Calcul du score global}
\label{subsec:sprint2-global-score}

Afin de disposer d'un indicateur unique pour comparer les expériences,
la plateforme calcule un \textbf{score global} à partir des métriques
sélectionnées par l'utilisateur.

\subsubsection{Normalisation des métriques}

Les métriques n'ont pas toutes la même échelle. Certaines, comme
l'Accuracy ou le F1-score, sont naturellement comprises entre 0 et 1.
D'autres, comme le MAE ou le RMSE, peuvent prendre des valeurs
beaucoup plus grandes sans borne supérieure claire.

Avant de calculer le score global, chaque métrique est donc
\textbf{normalisée dans l'intervalle $[0, 1]$} selon sa direction :

\begin{itemize}
    \item Pour les métriques de direction \textbf{max} (Accuracy,
    Precision, Recall, F1, $R^2$) dont la valeur est supérieure à 1
    mais inférieure ou égale à 100, la valeur est divisée par 100.
    Les valeurs déjà dans $[0, 1]$ sont conservées telles quelles.

    \item Pour les métriques de direction \textbf{min} (RMSE, MAE,
    MSE, loss), la transformation suivante est appliquée :
    \[
    score = \frac{1}{1 + max(valeur, 0)}
    \]
    Cette formule garantit que le résultat reste toujours dans
    $[0, 1]$, où une valeur proche de 0 est préférable.

    \item Les valeurs négatives sont ramenées à 0 et les valeurs
    supérieures à 100 à 1.
\end{itemize}

\subsubsection{Formule du score global}

Le calcul prend en compte l'ordre de priorité attribué aux trois
métriques principales :

\begin{itemize}
    \item \textbf{Priorité 1 : 50 \%} ;
    \item \textbf{Priorité 2 : 30 \%} ;
    \item \textbf{Priorité 3 : 20 \%}.
\end{itemize}

La formule utilisée est la suivante :

\[
\boxed{
Score_{global}
=
\frac{0{,}50 \times M_{1} + 0{,}30 \times M_{2} + 0{,}20 \times M_{3}}
{appliedWeight}
}
\]

où $M_{1}$, $M_{2}$ et $M_{3}$ représentent les métriques
normalisées selon leur priorité. Si une des métriques n'est pas
disponible pour un run donné, son poids n'est pas appliqué et le
dénominateur \texttt{appliedWeight} est ajusté en conséquence. Cette
renormalisation permet de conserver des scores comparables entre
des runs n'ayant pas tous les mêmes métriques.

Ainsi, la métrique considérée comme la plus importante influence
davantage le résultat, tandis que les deux autres restent prises en
compte dans l'évaluation globale.

Le score global facilite ensuite la comparaison entre plusieurs
expériences ou plusieurs versions d'un même modèle. Les mécanismes de
sélection, de promotion et de mise en production seront détaillés dans
le Sprint 3.


% =========================================================
% 4.6 MOTEUR D'ENTRAÎNEMENT
% =========================================================

\section{Mise en place du moteur d'entraînement}
\label{sec:sprint2-training-engine}

Afin de répondre à différents niveaux de maîtrise et à différents
besoins techniques, deux modes d'entraînement ont été intégrés à la
plateforme : un \textbf{mode standard}, basé sur un environnement déjà
préparé, et un \textbf{mode Docker personnalisé}, permettant à
l'utilisateur de définir son propre environnement.

% =========================================================
\subsection{Mode d'entraînement standard}
\label{subsec:sprint2-standard}

Le premier mode d'exécution est le \textbf{mode standard}. Il permet à
l'utilisateur de lancer rapidement un entraînement sans avoir à préparer
lui-même un environnement Docker.

Ce mode est principalement destiné aux premiers essais, aux prototypes,
aux tests rapides ainsi qu'aux utilisateurs qui souhaitent se concentrer
sur le script et les données plutôt que sur la préparation de
l'environnement technique.

Dans ce mode, la plateforme utilise une \textbf{image Docker fixe déjà
construite et prête à l'emploi}. Cette image contient notamment
\textit{NumPy}, \textit{Pandas}, \textit{Scikit-learn},
\textit{SciPy}, \textit{Joblib}, \textit{Matplotlib},
\textit{Seaborn}, \textit{XGBoost}, \textit{LightGBM},
\textit{PyTorch}, \textit{TorchVision}, \textit{OpenCV},
\textit{Pillow}, \textit{CloudPickle}, \textit{PyYAML},
\textit{Requests} et \textit{MLflow}.

L'un de ses principaux avantages est donc la rapidité : aucune nouvelle
image n'est construite avant chaque entraînement.

Cette simplicité impose cependant une limite. Le script Python doit
rester compatible avec les dépendances présentes dans l'image standard.

Une \textbf{analyse automatique du script} est donc réalisée avant
l'exécution. Le backend soumet le fichier à un service FastAPI dédié
qui analyse les imports via un parser AST et détecte les algorithmes
ML utilisés. Cette analyse permet de vérifier la compatibilité avec
le runtime standard, de déduire le type de problème
(classification ou régression) et le type de données
(tabulaire, textuel ou image).

Si les dépendances sont compatibles, le run peut être lancé. Dans le
cas contraire, l'utilisateur peut utiliser le mode Docker personnalisé.


% =========================================================
\subsubsection{Workflow du mode standard}
\label{subsubsec:sprint2-standard-workflow}

Après la validation du script, le backend envoie le fichier et le
dataset au \textbf{Pipeline Gateway}, un service FastAPI dédié qui
assure la passerelle entre le backend et Kubeflow Pipelines. Le
Gateway génère un fichier de configuration YAML contenant les
paramètres du run, puis soumet le pipeline au serveur Kubeflow.

Le pipeline standard est composé d'un seul step. Il utilise
l'image standard déjà construite, télécharge le dataset depuis le
Gateway puis exécute le script Python. Les métriques et les artefacts
sont automatiquement enregistrés dans MLflow.

Le diagramme de séquence \ref{fig:uml-standard-sequence} montre les
échanges entre les cinq composants impliqués dans un entraînement
standard : le frontend, le backend, le Pipeline Gateway, Kubeflow
et MLflow.

\begin{figure}[H]
    \centering

    \includegraphics[
        width=0.92\textwidth,
        keepaspectratio
    ]{images/uml_standard_sequence.png}

    \caption{Diagramme de séquence -- Entraînement standard}
    \label{fig:uml-standard-sequence}
\end{figure}

Le workflow peut être résumé comme suit :

\begin{center}

\fbox{
\parbox{7cm}{
\centering
\textbf{Script Python (.py) + Dataset}
}
}

\vspace{0.10cm}

$\downarrow$

\vspace{0.10cm}

\fbox{
\parbox{7cm}{
\centering
\textbf{Backend Spring Boot}\\
{\small Analyse AST et préparation du run}
}
}

\vspace{0.10cm}

$\downarrow$

\vspace{0.10cm}

\fbox{
\parbox{7cm}{
\centering
\textbf{Pipeline Gateway (FastAPI:8001)}\\
{\small Génération du YAML et soumission Kubeflow}
}
}

\vspace{0.10cm}

$\downarrow$

\vspace{0.10cm}

\fbox{
\parbox{7cm}{
\centering
\textbf{Kubeflow Pipeline (un seul step)}\\
{\small Image standard + téléchargement dataset + exécution}
}
}

\vspace{0.10cm}

$\downarrow$

\vspace{0.10cm}

\fbox{
\parbox{7cm}{
\centering
\textbf{MLflow}
$\rightarrow$
\textbf{Métriques, artefacts et résultats}
}
}

\end{center}

\begin{figure}[H]

    \centering

    \includegraphics[
        width=0.80\textwidth,
        keepaspectratio
    ]{images/sprint2_kubeflow_run.png}

    \caption{Exécution du pipeline d'entraînement standard dans Kubeflow}
    \label{fig:sprint2-kubeflow-run}

\end{figure}


% =========================================================
\subsection{Mode Docker personnalisé}
\label{subsec:sprint2-custom}

Le second mode d'exécution est le \textbf{mode Docker personnalisé}.
Il s'adresse principalement aux utilisateurs avancés qui souhaitent
garder la maîtrise de leur environnement d'entraînement.

Dans ce mode, l'utilisateur fournit le \textbf{script Python}, le
\textbf{dataset}, le fichier \textbf{requirements.txt} ainsi que le
\textbf{Dockerfile}.

Il peut ainsi utiliser des bibliothèques particulières, choisir leurs
versions et préparer un environnement adapté aux besoins de son modèle.

La plateforme reste cependant responsable de l'orchestration. Elle
modifie le Dockerfile fourni par l'utilisateur pour y injecter le
code platform (script d'orchestration MLflow, dépendances communes,
gestion des artefacts), prépare les fichiers nécessaires puis
construit l'image personnalisée.


% =========================================================
\subsubsection{Jobs Kubernetes de build et d'entraînement}
\label{subsubsec:sprint2-custom-jobs}

L'exécution repose sur deux Jobs Kubernetes successifs.

Le premier job utilise \textbf{Kaniko} pour construire l'image Docker
directement dans le cluster, sans dépendre d'un démon Docker. L'image
est ensuite poussée vers \textbf{Azure Container Registry} afin d'être
disponible pour le second job. Kaniko construit l'image à partir du
Dockerfile modifié et des fichiers préparés par la plateforme, puis
l'enregistre dans le registre avec un tag identifiant le projet et
l'expérience.

Le second job, nommé \texttt{mlops-custom-...}, utilise l'image créée
pour exécuter le script d'entraînement. Le dataset est monté via un
volume persistant partagé afin d'être accessible sans avoir été
intégré dans l'image.

Le diagramme de séquence \ref{fig:uml-custom-sequence} présente la
séquence en deux jobs : le build de l'image via Kaniko puis
l'exécution du script d'entraînement.

\begin{figure}[H]
    \centering

    \includegraphics[
        width=0.92\textwidth,
        keepaspectratio
    ]{images/uml_custom_sequence.png}

    \caption{Diagramme de séquence -- Entraînement Docker personnalisé}
    \label{fig:uml-custom-sequence}
\end{figure}

\begin{center}

\fbox{
\parbox{7.3cm}{
\centering
\textbf{Script + Dataset + Dockerfile + requirements.txt}
}
}

\vspace{0.10cm}

$\downarrow$

\vspace{0.10cm}

\fbox{
\parbox{7.3cm}{
\centering
\textbf{Backend Spring Boot}\\
{\small Injection du code platform dans le Dockerfile}
}
}

\vspace{0.10cm}

$\downarrow$

\vspace{0.10cm}

\fbox{
\parbox{7.3cm}{
\centering
\textbf{Job Kubernetes de build (Kaniko)}\\
{\small Construction et push vers Azure Container Registry}
}
}

\vspace{0.10cm}

$\downarrow$

\vspace{0.10cm}

\fbox{
\parbox{7.3cm}{
\centering
\textbf{Image Docker personnalisée (ACR)}
}
}

\vspace{0.10cm}

$\downarrow$

\vspace{0.10cm}

\fbox{
\parbox{7.3cm}{
\centering
\textbf{Job Kubernetes d'entraînement}\\
{\small Montage du dataset via PVC partagé, exécution du script}
}
}

\vspace{0.10cm}

$\downarrow$

\vspace{0.10cm}

\fbox{
\parbox{7.3cm}{
\centering
\textbf{MLflow}
$\rightarrow$
\textbf{Résultats}
}
}

\end{center}

La figure \ref{fig:sprint2-custom-jobs} présente un exemple des Jobs
créés lors d'un entraînement personnalisé.

\begin{figure}[H]

    \centering

    \includegraphics[
        width=0.76\textwidth,
        keepaspectratio
    ]{images/sprint2_custom_jobs.png}

    \caption{Jobs Kubernetes de construction de l'image et d'entraînement personnalisé}
    \label{fig:sprint2-custom-jobs}

\end{figure}

Cette séparation facilite également le diagnostic. Un échec du Job de
build indique généralement un problème lié à la préparation de
l'environnement (Dockerfile, dépendances, syntaxe), tandis qu'un
échec du Job d'entraînement oriente davantage l'analyse vers le
script, les données ou le modèle.


% =========================================================
\subsection{Comparaison des deux modes d'entraînement}
\label{subsec:sprint2-comparison}

Les deux modes d'entraînement restent complémentaires et répondent à
des besoins différents.

\begin{table}[H]

\centering
\small

\renewcommand{\arraystretch}{1.15}
\setlength{\tabcolsep}{4pt}

\begin{tabular}{|p{3.1cm}|p{5.1cm}|p{5.1cm}|}

\hline

\rowcolor{cyan!20}
\textbf{Critère}
&
\textbf{Mode standard}
&
\textbf{Mode Docker personnalisé}
\\
\hline

\textbf{Environnement}
&
Image fixe déjà préparée
&
Image propre au projet
\\
\hline

\rowcolor{cyan!7}
\textbf{Build avant entraînement}
&
Non
&
Oui (Kaniko)
\\
\hline

\textbf{Dépendances}
&
Disponibles dans l'image standard
&
Définies par l'utilisateur
\\
\hline

\rowcolor{cyan!7}
\textbf{Public visé}
&
Débutants, prototypes et tests rapides
&
Utilisateurs avancés et projets spécifiques
\\
\hline

\textbf{Avantage principal}
&
Simplicité et rapidité
&
Flexibilité et contrôle
\\
\hline

\end{tabular}

\caption{Comparaison des modes standard et Docker personnalisé}
\label{tab:sprint2-mode-comparison}

\end{table}


% =========================================================
% 4.7 SUIVI, NOTIFICATIONS ET RÉSULTATS
% =========================================================

\section{Suivi des expériences, notifications et résultats}
\label{sec:sprint2-tracking-results}


% =========================================================
\subsection{Suivi des expériences avec MLflow}
\label{subsec:sprint2-mlflow}

MLflow est intégré aux deux modes afin de conserver une trace des
expériences exécutées.

Pour chaque run, il permet de retrouver les métriques produites ainsi
que les artefacts associés. Le backend synchronise automatiquement
les informations entre MLflow et la base de données de la plateforme,
afin que les résultats soient consultables directement depuis
l'interface sans avoir à accéder à l'outil MLflow.

Cette intégration assure la traçabilité entre le projet, le run et
les résultats produits, tout en préparant les fonctions de
versionnement développées dans le sprint suivant.

\begin{figure}[H]

    \centering

    \includegraphics[
        width=0.80\textwidth,
        keepaspectratio
    ]{images/sprint2_mlflow_run.png}

    \caption{Suivi d'une expérience d'entraînement dans MLflow}
    \label{fig:sprint2-mlflow-run}

\end{figure}


% =========================================================
\subsection{Notifications liées aux runs}
\label{subsec:sprint2-notifications}

Les entraînements pouvant durer plusieurs minutes, un mécanisme de
notification en temps réel a été mis en place pour informer
l'utilisateur sans qu'il ait besoin de recharger la page.

Ce système repose sur \textbf{WebSocket} via le protocole STOMP
avec SockJS. Chaque utilisateur s'abonne à un topic dédié
(\texttt{/topic/users/\{userId\}/notifications}) et reçoit
directement les événements émis par le backend.

Lorsqu'un run se termine correctement, une notification confirmant la
réussite est envoyée. En cas d'échec, l'utilisateur est informé que
l'entraînement n'a pas abouti. Le même mécanisme est utilisé pour
d'autres événements du cycle de vie du modèle : démarrage et résultat
d'un déploiement, déclenchement d'un réentraînement ou encore alertes
de monitoring.

Côté administrateur, une diffusion séparée permet de suivre
l'ensemble des événements plateforme depuis un seul canal.

Les notifications opérationnelles ne passent pas par email. Ce canal
est réservé à la vérification de compte et à la réinitialisation du
mot de passe.

\begin{figure}[H]

    \centering

    \includegraphics[
        width=0.76\textwidth,
        keepaspectratio
    ]{images/sprint2_notifications.png}

    \caption{Notifications associées au succès ou à l'échec des runs}
    \label{fig:sprint2-notifications}

\end{figure}


% =========================================================
\subsection{Consultation des résultats}
\label{subsec:sprint2-results}

Les informations utiles sont finalement remontées vers le backend puis
affichées dans l'interface.

L'utilisateur peut consulter le statut du run, les métriques obtenues
ainsi que le score global calculé pour l'expérience sans devoir accéder
directement aux outils techniques.

Le score facilite également la comparaison entre plusieurs expériences
ou plusieurs versions d'un même modèle.

Une version satisfaisant les conditions de qualité peut ensuite devenir
candidate à la production. Les mécanismes de versionnement, de promotion
et de déploiement seront détaillés dans le Sprint 3.

\begin{figure}[H]

    \centering

    \includegraphics[
        width=0.80\textwidth,
        keepaspectratio
    ]{images/sprint2_training_results.png}

    \caption{Consultation des résultats d'entraînement depuis la plateforme}
    \label{fig:sprint2-training-results}

\end{figure}


% =========================================================
% 4.8 DIAGRAMMES DE SÉQUENCE UML
% =========================================================

\section{Diagrammes de séquence UML}
\label{sec:sprint2-uml}

Afin de mieux illustrer les interactions entre les différents
composants de la plateforme, deux diagrammes de séquence ont été
réalisés. Ils couvrent le mode d'entraînement standard et le mode
Docker personnalisé.


% =========================================================
% 4.9 BILAN DU SPRINT 2
% =========================================================

\section{Bilan du Sprint 2}
\label{sec:sprint2-bilan}

À la fin du Sprint 2, la plateforme disposait d'un premier parcours
réellement exploitable.

Les utilisateurs pouvaient s'authentifier, gérer leur espace, créer
leurs projets, importer les fichiers nécessaires et travailler avec
plusieurs types de données.

Les deux modes d'exécution étaient opérationnels : le mode standard
pour les tests rapides dans un environnement déjà préparé, et le mode
Docker personnalisé pour les projets nécessitant davantage de contrôle
sur les dépendances et l'environnement d'exécution.

La gestion des métriques permettait d'adapter l'évaluation aux problèmes
de classification ou de régression. Le score global apportait un
indicateur commun pour comparer les expériences tout en tenant compte
des priorités définies par l'utilisateur.

MLflow assurait le suivi des runs, tandis que les notifications WebSocket
et l'affichage des résultats facilitaient le suivi directement depuis
l'interface de la plateforme.


% =========================================================
% 4.10 CONCLUSION
% =========================================================

\section{Conclusion}
\label{sec:ch4-conclusion}

Ce deuxième sprint a permis de construire le cœur fonctionnel de la
plateforme et de relier la partie applicative au moteur d'entraînement
MLOps.

La solution permet désormais de gérer les utilisateurs et les projets,
de prendre en charge différents types de données, de configurer les
métriques, de choisir un mode d'exécution, de lancer un entraînement et
d'en suivre le résultat dans un environnement centralisé.

Le chapitre suivant sera consacré au Sprint 3, qui étend ce
fonctionnement au reste du cycle de vie du modèle, notamment au
versionnement, au déploiement, à la prédiction, au monitoring et au
réentraînement.

% =========================================================
% FIN DU CHAPITRE 4
% =========================================================
